\paragraph{谱零点的高阶 \texorpdfstring{$\ZZ_2$}{Z2} 轨道抵消}
在谱零点的单坐标判别、余类并结构与六倍窗 half-order dark band 已经建立之后，还可进一步把单个按位翻转对合的相消推广为一个多坐标 \(\ZZ_2\)-轨道抵消机制。其结果是：一旦某个频率 \(k\) 触发若干独立坐标上的 \(-1\) 相位翻转，则整个 \(\Omega_m\) 会按自由的 \((\ZZ/2\ZZ)^r\)-轨道分块，而每个轨道上的相位和都逐块归零。

\begin{theorem}[谱零点的高阶 \texorpdfstring{$\ZZ_2$}{Z2} 轨道抵消]\label{thm:xi-time-part70c-spectrum-zero-z2-orbit-cancellation}
\leanverified{paper\_xi\_time\_part70c\_spectrum\_zero\_z2\_orbit\_cancellation}
沿用定理 \ref{thm:fold-spectrum-zero-criterion} 与推论 \ref{cor:fold-spectrum-zero-involution} 的记号，记
\[
M:=F_{m+2},
\qquad
\widehat d_m(k)=\sum_{\omega\in\Omega_m}e^{-2\pi i kN(\omega)/M}.
\]
对任意 \(k\in \ZZ/(M\ZZ)\)，定义
\[
J(k):=
\left\{
j\in\{1,\dots,m\}:\ 2kF_{j+1}\equiv M\pmod{2M}
\right\},
\]
并令
\[
\Phi_k(\omega):=e^{-2\pi i kN(\omega)/M}
\qquad
(\omega\in \Omega_m).
\]
对每个 \(j\in J(k)\)，记
\[
\iota_j:\Omega_m\to \Omega_m,
\qquad
\iota_j(\omega):=\omega\oplus e_j.
\]
再定义
\[
G_k:=\langle \iota_j:\ j\in J(k)\rangle.
\]
则成立：
\begin{enumerate}
  \item \(G_k\cong (\ZZ/2\ZZ)^{|J(k)|}\)，且其在 \(\Omega_m\) 上的作用自由；
  \item 对任意
  \[
  g=\prod_{j\in S}\iota_j\in G_k,
  \]
  都有
  \[
  \Phi_k(g\omega)=(-1)^{|S|}\Phi_k(\omega)
  \qquad
  (\omega\in\Omega_m);
  \]
  \item 若 \(J(k)\neq \varnothing\)，则对任意 \(G_k\)-轨道 \(\mathcal O\subset \Omega_m\)，都有
  \[
  \sum_{\omega\in\mathcal O}\Phi_k(\omega)=0;
  \]
  特别地，
  \[
  \widehat d_m(k)=0
  \qquad\Longleftrightarrow\qquad
  J(k)\neq \varnothing.
  \]
\end{enumerate}
\end{theorem}
\begin{proof}
由定理 \ref{thm:fold-spectrum-zero-criterion}，
\[
\widehat d_m(k)=0
\qquad\Longleftrightarrow\qquad
\exists\,j\in\{1,\dots,m\}\ \text{使}\ 2kF_{j+1}\equiv M\pmod{2M},
\]
这正等价于 \(J(k)\neq \varnothing\)。

对每个 \(j\in J(k)\)，推论 \ref{cor:fold-spectrum-zero-involution} 已给出：\(\iota_j\) 是 \(\Omega_m\) 上无固定点的对合，且
\[
\Phi_k(\iota_j\omega)=-\Phi_k(\omega).
\]
若 \(j,j'\in J(k)\) 且 \(j\neq j'\)，则 \(\iota_j,\iota_{j'}\) 只翻转不同坐标，因此彼此交换，且各自平方为恒等。于是
\[
G_k\cong (\ZZ/2\ZZ)^{|J(k)|}.
\]
任一非平凡元素 \(g=\prod_{j\in S}\iota_j\) 都会翻转至少一个坐标，故不可能固定任何 \(\omega\in\Omega_m\)；从而 \(G_k\) 在 \(\Omega_m\) 上自由作用。

再对
\[
g=\prod_{j\in S}\iota_j
\]
重复使用 \(\Phi_k(\iota_j\omega)=-\Phi_k(\omega)\)，便得
\[
\Phi_k(g\omega)=(-1)^{|S|}\Phi_k(\omega).
\]
因此当 \(J(k)\neq \varnothing\) 时，对任意轨道
\[
\mathcal O=G_k\cdot \omega
\]
都有
\[
\sum_{\eta\in\mathcal O}\Phi_k(\eta)
=
\Phi_k(\omega)\sum_{S\subseteq J(k)}(-1)^{|S|}
=
\Phi_k(\omega)(1-1)^{|J(k)|}
=0.
\]
把 \(\Omega_m\) 分解为 \(G_k\)-轨道的并，即得
\[
\widehat d_m(k)=\sum_{\omega\in\Omega_m}\Phi_k(\omega)=0.
\]
与第一步的零点判别合并，即得最后的充要条件。证毕。
\end{proof}

\endinput
